%% capitulos/introducao.tex — Capítulo 1: Introdução
%%
%% Bloco escrito por Claude com base na qualificação de 2025 e em buscas
%% bibliográficas no scite (DOIs verificados). Conteúdo revisado e aprovado;
%% a marcação de rascunho foi removida em 2026-09-29.

\chapter{Introdução}
\label{ch:introducao}

% WRITER:BEGIN:introducao/full

\section{Contexto da análise de imagens médicas no diagnóstico clínico}
\label{sec:intro-contexto}

A imagem médica consolidou-se como um dos pilares do diagnóstico clínico
contemporâneo. Modalidades como tomografia computadorizada, ressonância
magnética, radiografia, ultrassonografia, tomografia por emissão de pósitrons
e a histopatologia digital fornecem representações detalhadas de estruturas
internas e teciduais que sustentam decisões de triagem, estadiamento e
acompanhamento terapêutico em praticamente todas as especialidades médicas
\cite{jiang2023}. O volume de dados gerado por esses exames cresce de modo
contínuo, ao mesmo tempo em que sua interpretação permanece uma tarefa que
exige treinamento especializado e está sujeita a variabilidade inter e
intraobservador. Esse cenário motiva o desenvolvimento de ferramentas
computacionais capazes de oferecer suporte quantitativo à análise dessas
imagens, notadamente em tarefas de classificação, segmentação e detecção
de anomalias.

\section{Avanços em aprendizado profundo aplicados à medicina}
\label{sec:intro-dl}

Nas últimas décadas o aprendizado profundo (\emph{deep learning}, DL)
emergiu como o paradigma dominante para o processamento de imagens médicas.
Em particular, redes neurais convolucionais profundas, exemplificadas pela
família de redes residuais introduzida por \citeonline{he2016}, alcançaram
desempenho próximo ou superior ao de especialistas humanos em diversas
tarefas de classificação e detecção \cite{jiang2023}. Esse progresso é
sustentado, de um lado, pelo aumento sistemático do poder computacional e,
de outro, pela disponibilização de coleções públicas padronizadas de imagens
médicas que viabilizam treinamento, comparação e reprodutibilidade de
modelos. Entre essas coleções destaca-se o \emph{benchmark} MedMNIST v2
\cite{medmnist2023}, que reúne dezoito conjuntos de classificação 2D e 3D
oriundos de patologia, radiologia e oftalmologia, padronizados em baixa
resolução para permitir prototipagem rápida sem perder representatividade
das modalidades clínicas.

Apesar desses avanços, modelos profundos operam de modo essencialmente
opaco do ponto de vista clínico: a relação entre os parâmetros aprendidos
e as decisões de saída não é trivialmente legível, configurando o problema
classicamente designado por \emph{caixa-preta}. Levantamentos recentes
apontam que essa limitação constitui hoje a principal barreira à adoção
de DL em fluxos diagnósticos sensíveis, mais relevante até que o ganho
marginal de acurácia em métricas tradicionais \cite{teng2022}. Nesse
contexto, a interpretabilidade dos modelos depende não apenas de explicações
\emph{post hoc} sobre saídas individuais, mas também de instrumentos
quantitativos que caracterizem o estado interno da rede durante o
treinamento.

\section{Justificativa do uso de medidas de complexidade}
\label{sec:intro-justificativa}

Métricas tradicionais de avaliação --- acurácia, perda de validação,
área sob a curva ROC --- capturam o comportamento do modelo sobre um
conjunto de dados específico, mas não oferecem leitura direta sobre a
organização estatística dos pesos aprendidos. Essa lacuna é particularmente
visível no fenômeno do \emph{overfitting}: redes com elevada capacidade
podem memorizar o conjunto de treinamento de modo a atingir acurácia
quase perfeita sobre ele, ao mesmo tempo em que generalizam mal para
dados novos. Análises teóricas recentes da dinâmica de generalização em
regimes de alta dimensionalidade indicam que os pesos atravessam fases
qualitativamente distintas durante o treinamento, com implicações diretas
sobre a capacidade preditiva do modelo \cite{advani2020}, e métodos
dedicados à detecção de \emph{overfitting} em aplicações clínicas têm
sido propostos como resposta prática a esse problema \cite{zhang2025}.

A presente dissertação parte da hipótese de que medidas de complexidade
aplicadas diretamente aos pesos da rede podem servir como marcadores
objetivos da maturação do treinamento e do início do sobreajuste. Duas
famílias de medidas dão sustentação a essa abordagem. A primeira é a
Entropia de Amostra (\emph{Sample Entropy}, SampEn) introduzida por
\citeonline{richman2000} como refinamento da Entropia Aproximada de
\citeonline{pincus1991}; SampEn quantifica, sem viés sistemático associado a
contagens de auto-coincidência, a regularidade de uma sequência de valores
e tem aplicação consolidada na análise de sinais fisiológicos. Sua extensão
bidimensional, SampEn2D, foi formalizada por \citeonline{silva2016} para
caracterizar a irregularidade de texturas espaciais em imagens biomédicas,
e é diretamente aplicável à matriz de pesos de uma camada convolucional.
A análise multiescala (Multiscale Entropy, MSE) proposta por
\citeonline{costa2002} complementa SampEn ao revelar regularidades que se
manifestam em diferentes resoluções da mesma sequência ou matriz.

A segunda família é a complexidade estatística de López-Ruiz, Mancini e
Calbet \cite{lmc1995}, definida pelo produto entre uma medida de
incerteza, tipicamente a entropia de Shannon, e uma medida de
desequilíbrio em relação à distribuição uniforme. Diferentemente de
medidas puramente entrópicas, a complexidade LMC anula-se nos extremos
de ordem perfeita e desordem perfeita, atingindo valores máximos em
regimes intermediários de auto-organização \cite{anteneodo1996}.
Aplicada à distribuição dos pesos de uma rede neural, espera-se que ela
descreva quantitativamente a transição entre o estado inicial aleatório
e o estado treinado. A combinação entre SampEn (sensível à estrutura
sequencial e espacial dos pesos) e LMC (sensível à forma da distribuição
marginal) constitui o eixo metodológico desta investigação.

\subsection{Justificativa prática: o ponto de vista de quem treina a rede}
\label{sec:intro-just-dev}

A justificativa apresentada até aqui é de natureza metodológica: as medidas de
complexidade oferecem uma leitura dos pesos que as métricas de desempenho não
oferecem. Há, porém, uma segunda justificativa, de natureza prática, que motiva
este trabalho de modo mais imediato e é formulada deliberadamente do ponto de
vista de quem desenvolve e treina a rede.

\subsubsection*{A dependência atual dos dados rotulados}

Para saber se uma rede aprendeu ou apenas memorizou, o desenvolvedor depende
hoje, integralmente, dos dados. É preciso construir conjuntos de treino,
validação e teste bem particionados, corretamente rotulados e representativos da
população de interesse. Cada um desses requisitos custa tempo e recursos, e em
imagens médicas frequentemente esbarra na indisponibilidade de especialistas
para a anotação. O diagnóstico do treinamento é, portanto, função do orçamento
de anotação disponível --- e não uma propriedade que se possa ler diretamente do
modelo.

O procedimento padrão para decidir quando interromper o ajuste, a parada
antecipada, monitora o erro em um conjunto de validação separado
\cite{prechelt1998}. O procedimento é sólido, mas transfere ao conjunto de
validação toda a responsabilidade pelo diagnóstico: se ele é pequeno,
desbalanceado ou pouco representativo, a estimativa do momento de parada
torna-se ruidosa, e o desenvolvedor não dispõe de sinal alternativo com que
confrontá-la. Em outras palavras, justamente no regime de poucos dados ---
aquele em que o sobreajuste é mais provável --- o instrumento de medida é o
menos confiável.

\subsubsection*{O que muda com um indicador calculado sobre os pesos}

As medidas investigadas nesta dissertação são calculadas exclusivamente a partir
dos pesos. Não consomem imagem, não consomem rótulo e não dependem de nenhuma
partição. Caso sejam capazes de sinalizar a transição para o regime de
memorização, o desenvolvedor passa a dispor de uma segunda leitura do
treinamento, obtida por caminho independente daquele do conjunto de validação.
As duas leituras podem então ser confrontadas: quando concordam, a decisão de
parada ganha confiança; quando discordam, há indício de que a validação não é
representativa.

Note-se que a proposta não é substituir a parada antecipada clássica, cuja
fundamentação é sólida, mas acrescentar-lhe uma fonte de evidência independente.

\subsubsection*{Casos de uso}

\begin{enumerate}
    \item \textbf{Decidir quando interromper o treinamento.} Um aviso emitido
    antes da degradação observável na validação amplia a janela de decisão, em
    lugar de constatar a perda depois de consumada.

    \item \textbf{Economizar recursos computacionais.} Detectar que a rede
    deixou de aprender evita épocas de treinamento que não trarão ganho, com
    impacto direto no custo de processamento.

    \item \textbf{Comparar arquiteturas sem ampliar a anotação.} Como a medida
    não consome dados rotulados, a comparação entre arquiteturas candidatas não
    exige a ampliação do conjunto de validação.

    \item \textbf{Diagnosticar quando a validação é pequena.} Em cenários nos
    quais o conjunto de validação é pequeno demais para ser confiável, um
    indicador independente fornece evidência adicional --- precisamente o cenário
    típico em imagens médicas.
\end{enumerate}

\subsubsection*{Entrega prática e delimitação}

A contribuição não se esgota na constatação empírica. Entrega-se um monitor
implementado em PyTorch, acoplável ao laço de treinamento de terceiros, que
calcula os indicadores a cada época e devolve um estado do aprendizado, conforme
descrito na Seção~\ref{sec:monitor}. O custo computacional é baixo quando a
medida incide sobre a camada densa.

Registre-se, por honestidade, que a utilidade prática aqui descrita é
\emph{condicional} ao sucesso empírico dos indicadores, e que os resultados do
Capítulo~\ref{ch:resultados} mostram que esse sucesso não é uniforme: depende da
medida empregada. As implicações dessa delimitação para o uso prático são
discutidas no Capítulo~\ref{ch:conclusao}.

\section{Hipótese, objetivos e escopo}
\label{sec:intro-objetivos}

A hipótese central deste trabalho é que a complexidade estatística dos
pesos de uma rede neural profunda evolui de modo sistemático durante
o treinamento e constitui um marcador objetivo do início do
\emph{overfitting}. Mais especificamente, postula-se que a trajetória
das medidas LMC e SampEn ao longo dos passos de otimização exibe
inflexão não monotônica, e que essa inflexão precede ou coincide com a
saturação da acurácia de validação.

A pesquisa adota um \textbf{escopo deliberadamente restrito} a fim de
permitir o controle das variáveis de interesse e a reprodutibilidade
plena dos experimentos:

\begin{itemize}
  \item \textbf{Arquitetura:} ResNet-18 \cite{he2016}, escolhida
        por sua importância histórica, profundidade moderada e papel de
        \emph{baseline} consagrado em classificação de imagens médicas;
        parte dos experimentos é repetida com a DenseNet-121 para
        verificar a dependência em relação à arquitetura.
  \item \textbf{Conjunto de dados:} MedNIST, com seis classes de imagens
        médicas de diferentes modalidades (Seção~\ref{sec:dados}).
  \item \textbf{Pesos analisados:} apenas os da camada densa
        (classificadora), calculando-se as medidas sobre a camada inteira.
  \item \textbf{Tarefa única:} classificação multiclasse, excluindo-se
        explicitamente segmentação semântica, detecção e tarefas
        ordinais.
\end{itemize}

Arquiteturas alternativas como U-Net, conjuntos de dados de maior escala
como NIH ChestX-ray14 e TCIA, e tarefas adicionais como segmentação ---
mencionados no projeto de qualificação --- ficam relegados a trabalhos
futuros (\autoref{ch:conclusao}).

Os objetivos específicos são:

\begin{enumerate}
  \item Implementar SampEn1D, SampEn2D e a complexidade LMC com
        validação numérica em casos analiticamente conhecidos
        (distribuições uniforme, gaussiana e séries periódicas), conforme
        o \autoref{ap:validacao-numerica}.
  \item Extrair os pesos da camada densa a cada época, durante o
        treinamento, excluindo o vetor de viés da análise.
  \item Calcular a trajetória das medidas de complexidade em execuções
        com corrupção de rótulos e em seus controles pareados, e
        confrontá-la com a evolução da perda e da acurácia de validação.
  \item Avaliar se a inflexão das trajetórias LMC e SampEn coincide,
        antecipa ou sucede o início observável do sobreajuste, e
        comparar os comportamentos entre níveis de corrupção e entre
        arquiteturas.
\end{enumerate}

\section{Organização do trabalho}
\label{sec:intro-organizacao}

O restante desta dissertação está organizado em quatro capítulos. O
\autoref{ch:fundamentacao} apresenta a fundamentação teórica das
medidas de complexidade utilizadas (Entropia de Shannon, ApEn, SampEn,
SampEn2D, complexidade LMC) e situa-as no contexto da literatura
de redes neurais profundas e análise de séries de pesos. O
\autoref{ch:metodologia} descreve o conjunto de dados, as arquiteturas,
o protocolo de treinamento, a indução de memorização por corrupção de
rótulos, a camada densa como objeto de medida e os parâmetros adotados
para cada medida ($m$, $r$, número de \emph{bins}). O
\autoref{ch:resultados} apresenta as trajetórias de complexidade da
camada densa nas execuções com corrupção de rótulos e em seus controles
pareados, e sua relação com a perda e a acurácia de validação. O
\autoref{ch:conclusao} sintetiza as evidências obtidas em relação à
hipótese central, discute as limitações do recorte adotado e indica
direções para extensões futuras envolvendo arquiteturas, modalidades
de imagem e tarefas adicionais.

% WRITER:END:introducao/full
